\begin{problem}[2019年第36届全国中学生物理竞赛复赛][雷电云与闪电物理]{110}
闪电是地球上最壮丽的自然现象之一。人们对闪电进行了大量研究，近年来还观测到闪电导致的瞬间发光和伽玛射线暴等新现象。闪电通常由雷电云（离地 $6\sim12\,\mathrm{km}$）放电产生，多数闪电发生在云内，少数到达地面。由于云内冰状颗粒相互碰撞，小颗粒冰晶带正电，随气流上浮到云上端；较大颗粒带负电，下坠到云底端（见图a）。云中闪电中和了云内的正负电荷，而云地闪电则把负电荷释放到地面。
\begin{subproblem}
利用高空气球携带的电场测量仪测量高空中某圆柱形空域雷电云内的电场，其强度可视为均匀分布，大小为 $0.15\,\mathrm{MV/m}$。该圆柱区域的中轴线垂直于地面，半径为 $2.5\,\mathrm{km}$，高度为 $1.0\,\mathrm{km}$。求该区域上下两端的电势差、正电荷总量以及携带的总电能。已知真空介电常量 $\varepsilon_{0}=8.85\times10^{-12}\,\mathrm{F/m}$。
\end{subproblem}
\begin{subproblem}
在起电过程中，雷电云上下两端电荷会随时间指数增加。当地表电场大于 $1.0\,\mathrm{kV/m}$ 时，就会发生云地闪电，因此地表电场很少超过 $10\,\mathrm{kV/m}$。假定1）中所述的雷电云从高空缓慢整体下移，直至其负电荷层离地高度为 $6.0\,\mathrm{km}$ 时暂时保持稳定，地面为良导体，试估算此雷电云正下方产生的地表电场强度。
\end{subproblem}
\begin{subproblem}
云地闪电通常由带电云底端带负电的冰晶颗粒尖端放电触发，先形成一条指向地面的放电细路径（直径为厘米量级），该细路径随时间向下延伸，并导致周围空气不断电离，逐渐形成以原细路径（横截面大小可视为不变）为轴的粗圆柱形带电体，最后接近地面形成云地闪电通道。该闪电通道垂直于地面，所带负电荷总量为 $2.5\,\mathrm{C}$（原细放电路径内所带电量相对很小），闪电通道（中心放电细路径除外）内部电场强度大小相等。假设闪电通道的长度远大于其直径，闪电通道的直径远大于中心放电细路径的直径，且在闪电通道连通云地前的极短时间内，闪电通道内部的电荷分布可视为稳定分布。已知大气的电场击穿阈值为 $3.0\,\mathrm{MV/m}$，试估算该云地闪电通道的直径，并导出闪电通道（中心放电细路径除外）内的电荷密度径向分布的表达式。
\end{subproblem}
\begin{subproblem}
闪电通道连通云地后，云底和通道内部的负电荷迅速流向地面；闪电区域的温度骤然上升到数万摄氏度，导致其中的空气电离，形成等离子体，放出强光，同时通道会剧烈膨胀，产生雷声，闪电的放电电流经过约 $10\,\mu\mathrm{s}$ 时间即可达数万安培。在通道底部（接近地面）向四周辐射出频率约为 $30\,\mathrm{kHz}$ 的很强的无线电波。由于频率低于 $20\,\mathrm{MHz}$（此即所谓电离层截止频率）的电磁波不能进入电离层内部，该无线电波会加热电离层底部（离地约 $80\,\mathrm{km}$）的等离子体，闪电电流一旦超过某阈值将导致该电离层底部瞬间发光，形成一个以强无线电波波源（通道底部）正上方对应的电离层底部为中心的光环，最大直径可延伸到数百公里。试画出电离层底部光环产生与扩展的物理过程示意图，并计算光环半径为 $100\,\mathrm{km}$ 时光环扩张的径向速度。
\end{subproblem}
\begin{subproblem}
球形闪电（球闪）的微波空泡模型认为球闪是一个球形等离子体微波空腔（空泡）。当闪电微波较弱时，不足以形成微波空泡，会向太空辐射，穿透电离层，可被卫星观测到。实际上，卫星确实观测到了这种微波辐射。但卫星观测信号易受电离层色散的干扰，携带探测器的高空气球可到达雷电云上方观测，以避免此类干扰。为了在离地 $12\,\mathrm{km}$ 的高空观测闪电发出的微波信号，需要在该区域悬浮一个载荷（包括气球材料和探测器）为 $50\,\mathrm{kg}$ 的高空氦气球，求此气球在高空该区域悬浮时的体积。已知在离地 $12\,\mathrm{km}$ 的高度以下，大气温度随高度每升高 $1\,\mathrm{km}$ 下降 $5.0\,\mathrm{K}$，地面温度 $T_{0}=290\,\mathrm{K}$，地面压强 $p_{0}=1.01\times10^{5}\,\mathrm{Pa}$，空气摩尔质量 $M=29\,\mathrm{g/mol}$；气球内氦气密度（在离地高度 $12\,\mathrm{km}$ 处的值） $\rho_{\mathrm{He}}=0.18\,\mathrm{kg/m^{3}}$。重力加速度 $g=9.8\,\mathrm{m/s^{2}}$，气体普适常量 $R=8.31\,\mathrm{J/(K\cdot mol)}$。
\end{subproblem}
\end{problem}